% resume-entry-inline-optional.tex — optional entry fields under `entrymeta=inline'.
%
% Purpose (issue #230)
%   Rule 5 of AGENTS.md on the horizontal axis. `entrymeta=inline' introduces
%   the first separator that falls *inside* a line of the entry heading, and a
%   separator emitted unconditionally would put a `|' at the start of a line, or
%   at the end of one, whenever a cell it was meant to join is absent.
%
%   The column form cannot have this defect, because `\hfill' is not a
%   separator: with an empty left cell it simply puts the surviving cell at the
%   right margin. So the entire optional-field question is new under `inline',
%   and it is answered here rather than only in the regression suite because a
%   .tlg records that the separator *code* was not reached, while this records
%   that no `|' reached the page.
%
%   Every shape the four fields can take, minus the full one (which
%   resume-entry-inline-order.tex covers):
%
%     1. title + organization + location, no dates
%     2. title + dates, no organization or location
%     3. title + location, no organization  — the stray-separator shape: nothing
%        precedes the location on its line, so it must open flush left with no
%        mark before it
%     4. title + organization, no location or dates
%     5. title alone — no second heading line at all
%
%   The assertion is the whole baseline, but the property to read it for is
%   simple: no line begins with `|', no line ends with `|', and no `|' appears
%   with nothing on one side of it.
%
%   The committed baseline (resume-entry-inline-optional.expected.txt) is the
%   assertion; regenerate it only for an intended, reviewed output change.
%
% Compile: lualatex -halt-on-error -interaction=nonstopmode resume-entry-inline-optional.tex
% (run from the repository root, or with the root on TEXINPUTS as run.sh does)
\documentclass[entrymeta=inline]{careerdossier-resume}
\CDossierSetup{
  name     = {Ada Lovelace},
  headline = {Accessibility Engineer},
  email    = {ada@example.com},
}
\begin{document}
\MakeCDossierHeader
\CDossierSection{Experience}
\begin{CDossierEntry}[
  title        = {Dateless Engineer},
  organization = {Example Labs},
  location     = {Toronto, ON}
]
  \begin{CDossierItemize}
    \item An entry whose dates are absent keeps its title line unseparated.
  \end{CDossierItemize}
\end{CDossierEntry}
\begin{CDossierEntry}[
  title = {Independent Consultant},
  dates = {2021--2023}
]
  \begin{CDossierItemize}
    \item An entry with dates and no second heading line.
  \end{CDossierItemize}
\end{CDossierEntry}
\begin{CDossierEntry}[
  title    = {Volunteer Reviewer},
  location = {Remote}
]
  \begin{CDossierItemize}
    \item A location with no organization opens its line with no separator.
  \end{CDossierItemize}
\end{CDossierEntry}
\begin{CDossierEntry}[
  title        = {Workshop Facilitator},
  organization = {Example Trust}
]
  \begin{CDossierItemize}
    \item An organization with no location ends its line with no separator.
  \end{CDossierItemize}
\end{CDossierEntry}
\begin{CDossierEntry}[
  title = {Title Only}
]
  \begin{CDossierItemize}
    \item A title alone emits no second heading line and no separator.
  \end{CDossierItemize}
\end{CDossierEntry}
\end{document}
